% answers_short.tex —— 2015 年·上海卷·文科 · 简短答案速查
% 由 tools/make_short_answers.py 自 answers.tex 的 \gksolhead/\gkanswer 提取，
% 并与 2015/2015_数学真题_上海文_答案要点.md 逐题比对后生成，勿手改（改源 MD 后重跑）。

\begin{gknotice}
\textbf{使用说明}：本节只给最终答案便于快速核对，完整推导见解析版。答案为编写者独立推导，\textbf{非官方标准答案}。
\end{gknotice}

\gkpart{填空题}
\begin{itemize}[leftmargin=2.6em]
  \item[\textbf{1.}] \(\pi\)
  \item[\textbf{2.}] \(\{1,4\}\)
  \item[\textbf{3.}] \(\dfrac{1}{4}+\dfrac{1}{2}\mathrm{i}\)
  \item[\textbf{4.}] \(-\dfrac{2}{3}\)
  \item[\textbf{5.}] \(16\)
  \item[\textbf{6.}] \(4\)
  \item[\textbf{7.}] \(2\)
  \item[\textbf{8.}] \(x=2\)
  \item[\textbf{9.}] \(3\)
  \item[\textbf{10.}] \(120\)
  \item[\textbf{11.}] \(240\)
  \item[\textbf{12.}] \(\dfrac{x^{2}}{4}-\dfrac{y^{2}}{4}=1\)
  \item[\textbf{13.}] \(3+\sqrt{5}\)
  \item[\textbf{14.}] \(8\)
\end{itemize}

\gkpart{选择题}
\begin{center}\begin{tabular}{ccccc}
\toprule
题号 & 15 & 16 & 17 & 18 \\
\midrule
答案 & A & B & D & A \\
\bottomrule
\end{tabular}\end{center}

\gkpart{解答题}
\begin{itemize}[leftmargin=2.6em]
  \item[\textbf{19.}] \(V_{P-AOC}=\dfrac{1}{3}\)；\(\cos\theta=\dfrac{\sqrt{10}}{10}\)。
  \item[\textbf{20.}] (1) \(a=0\) 时 \(f(x)\) 为奇函数；\(a\neq0\) 时 \(f(x)\) 为非奇非偶函数；(2) \(a\in(1,3)\) 时 \(f(x)\) 在 \([1,2]\) 上严格递增
  \item[\textbf{21.}] (1) \(t_{1}=\dfrac{3}{8}\) 小时；\(f(t_{1})=\dfrac{3\sqrt{41}}{8}\) 千米；(2) \(f(t)=\sqrt{25t^{2}-42t+18}\)（\(\dfrac{3}{8}\leqslant t\leqslant\dfrac{7}{8}\)）；\(f(t)\) 在 \(\left[\dfrac{3}{8},\dfrac{7}{8}\right]\) 上的最大值为 \(\dfrac{3\sqrt{41}}{8}\approx 2.40<3\)，故不超过 3
  \item[\textbf{22.}] (1) \(d=\dfrac{|x_{1}y_{2}-x_{2}y_{1}|}{\sqrt{x_{1}^{2}+y_{1}^{2}}}\)；且 \(S=\dfrac{1}{2}|x_{1}y_{2}-x_{2}y_{1}|\)；(2) \(k=-1\) 或 \(k=-\dfrac{1}{5}\)；(3) \(m=-\dfrac{1}{2}\)（此时恒有 \(S=\dfrac{\sqrt{2}}{4}\)，与 \(l_{1}\)、\(l_{2}\) 的变动无关）
  \item[\textbf{23.}] (1) \(a_{n}=6n-5\)；(2) 成立（\(b_{n}\leqslant b_{n_{0}}\) 对一切 \(n\in\mathbf{N}^{\ast}\)）；(3) \(\lambda\in\left(-\dfrac{1}{4},0\right)\)
\end{itemize}
